\documentclass[11pt,a4paper]{article}
\usepackage[utf8]{inputenc}
\usepackage[T1]{fontenc}
\usepackage{amsmath,amssymb}
\usepackage[margin=25mm]{geometry}
\title{Student Exam Tools: Visuell verstehen}
\author{Studentisch entwickelt. Gemeinsam besser lernen.}
\date{}
\begin{document}
\maketitle
\section*{1. Gleiche Teile (Klasse 5--6)}
Berechne $\frac34+\frac16$ mit einem gemeinsamen Nenner. Erkläre, warum gleich große Teilstücke nötig sind.
\section*{2. Funktion und Probe (Schule)}
Bestimme für $f(x)=x^2-4x+3$ die reellen Nullstellen und den Scheitelpunkt. Prüfe die Nullstellen im ursprünglichen Term; beschreibe den Graphen ohne reine Farbverweise.
\section*{3. Komplexe Wurzeln (Uni-Einstieg)}
Bestimme alle $z\in\mathbb C$ mit $z^3=8i$. Gib den Betrag und die drei Winkel an. Prüfe jede Lösung.
\section*{4. Matrixwirkung (Uni-Einstieg)}
Es sei $A=\begin{pmatrix}2&1\\1&2\end{pmatrix}$. Berechne $Av$ für $v=(1,1)^\mathsf T$ und $v=(1,-1)^\mathsf T$. Deute die Streckung.
\bigskip
Die Lösungen sind absichtlich getrennt. Sieh erst nach einem eigenen Versuch in die durchgerechneten Beispiele im Repository.
\end{document}
